\chapter{Taking opioid medicines safely}
\index{Taking opioid medicines safely}

\medskip
\noindent\textit{Educational reference; always seek professional advice for individual care.}

\section*{Definition and Overview}
Opioids are a group of powerful pain-relieving medicines that work by binding to receptors in the brain and spinal cord. Common examples in many countries include codeine, tramadol, oxycodone, morphine, hydromorphone, tapentadol, and fentanyl patches. They are among the most effective treatments available for moderate to severe pain, and they are essential in hospital, after surgery, and for people with cancer or at the end of life. However, they are also potentially dangerous medicines, with a real risk of dependence, overdose, and harm if used incorrectly or for longer than needed. This chapter is about using them safely: getting the benefit without the harm.

The key principles are simple. Opioids should be used for a clear reason, at the lowest effective dose, for the shortest useful time, and under the supervision of a doctor. For most acute pain, such as pain after an injury or operation, a short course of a few days works well, and the dose is then reduced and stopped. For long-term pain, opioids are now prescribed much more cautiously than in the past, because the evidence shows that their benefit falls over time while the risks of dependence and other harms rise. In many countries, access to opioids is controlled, and codeine is available only with a prescription, reflecting a national effort to reduce opioid harm.

This chapter explains how opioids work, the situations in which they are appropriate, the side effects and risks to watch for, and the practical habits that keep people safe, including safe storage, safe disposal, and what to do if an overdose is suspected. It is written for patients, families, and students as a general reference.

\section*{Epidemiology}
Opioids are among the most prescribed medicines in many countries, with millions of prescriptions dispensed each year. Codeine and oxycodone have historically been the most commonly used, and although overall prescribing has fallen since the rescheduling of codeine to prescription-only in 2018, opioids are still widely used for acute and chronic pain. Opioid use is more common in older age groups, who have more painful conditions, and in people with chronic musculoskeletal pain, back pain, and cancer. Regional and rural some countries have higher rates of use than the cities.

The harm from opioids has been a major public health concern. Over the two decades to the mid-2010s, opioid-related deaths in many countries rose substantially, driven partly by prescription opioids and partly by diverted or illicit supplies. Since then, tighter regulation, safer prescribing, and greater awareness have stabilised and in some groups reduced the death rate, but opioids still account for a significant proportion of drug-induced deaths in many countries, the majority of which involve other sedating substances such as alcohol and benzodiazepines taken at the same time. Understanding these numbers helps explain why doctors now ask searching questions before prescribing and why "just a script" is never the whole answer.

\section*{Aetiology and Pathophysiology}
Opioids produce pain relief by binding to specific receptors, principally the mu-opioid receptor, which are spread throughout the brain, spinal cord, and gut. When activated, these receptors reduce the transmission of pain signals up the spinal cord and dampen the brain's perception of pain. They also produce a feeling of well-being or euphoria, which is part of why they work so well and part of why they are so easily misused. The same receptors, when strongly activated, slow breathing, constipate the bowel, and suppress coughing.

Dependence develops because the brain adapts to the drug. With repeated use, the brain adjusts its own chemistry to compensate, so that stopping the medicine suddenly causes a withdrawal syndrome of anxiety, sweating, cramps, and a racing heart. Tolerance means the same dose produces less effect over time, tempting people to take more. Overdose occurs when too much opioid overwhelms the brain's breathing centre: breathing becomes slow, shallow, or stops, leading to loss of consciousness and, without rapid help, death. The risk is greatly magnified when opioids are combined with alcohol or other sedating medicines.

\begin{itemize}
    \item \textbf{Pain relief:} the intended use, for moderate to severe pain after surgery, injury, or serious illness
    \item \textbf{Dependence:} the brain adapts to repeated use, so stopping suddenly causes withdrawal
    \item \textbf{Tolerance:} the same dose works less well over time, leading to dose escalation
    \item \textbf{Respiratory depression:} the dangerous effect that causes overdose, worsened by alcohol and other sedatives
    \item \textbf{Constipation:} opioids slow the gut, and laxatives are often needed alongside
    \item \textbf{Diversion:} unused tablets can be taken by others, so secure storage and safe disposal are essential
\end{itemize}

\section*{Clinical Features}
This section covers the effects a person can expect when taking opioids properly, the side effects to report, and the danger signs.

\begin{itemize}
    \item \textbf{Pain relief:} the medicine should noticeably reduce pain, allowing normal activity and sleep
    \item \textbf{Sedation:} drowsiness is common in the first days and with higher doses, and it impairs driving and operating machinery
    \item \textbf{Constipation:} nearly universal with regular opioid use; it should be prevented with fluids, fibre, and laxatives
    \item \textbf{Nausea:} common when starting, and usually settles, but should be reported if it persists
    \item \textbf{Dizziness and light-headedness:} can occur, especially when standing up quickly, and increases fall risk in older people
    \item \textbf{Itching and dry mouth:} common minor effects, usually manageable
    \item \textbf{Dangerous drowsiness:} severe sleepiness, slow or irregular breathing, or being hard to wake are overdose signs and are emergencies
    \item \textbf{Withdrawal:} if a dose is missed or the medicine is stopped suddenly, anxiety, sweating, cramps, and restlessness can occur
\end{itemize}

\section*{Differential Diagnosis}
The important diagnostic question in safe prescribing is whether an opioid is the right medicine at all, and what alternatives might be safer.

\begin{itemize}
    \item \textbf{Paracetamol and NSAIDs:} the first-line medicines for most mild to moderate pain, with far lower risk
    \item \textbf{Non-medication treatment:} physiotherapy, exercise, heat and cold, and psychological approaches are effective for many pains
    \item \textbf{Nerve pain medicines:} gabapentin, pregabalin, and some antidepressants are better choices for neuropathic (nerve) pain
    \item \textbf{Local anaesthetics and nerve blocks:} can provide targeted pain relief after surgery
    \item \textbf{Cancer and end-of-life pain:} opioids remain the mainstay, often combined with other strategies
    \item \textbf{Opioid use disorder:} when use has become compulsive despite harm, specialist treatment is needed rather than a simple prescription
\end{itemize}

\section*{When to Seek Medical Care}
Anyone prescribed an opioid should understand what to report to their doctor: pain that is not controlled, side effects that interfere with life, or a need for the medicine beyond the planned duration. For chronic pain, the GP will review the dose regularly and support a plan for tapering or for using other treatments. If you are caring for someone taking opioids, know the danger signs and do not hesitate to act on them.

\textbf{Contact your local emergency services for:}
\begin{itemize}
    \item Severe drowsiness or being difficult to wake
    \item Slow, shallow, or irregular breathing
    \item Blue or grey lips or skin, or cold, clammy skin
    \item Pinpoint pupils with unresponsiveness, which suggest opioid overdose
    \item Convulsions or collapse
\end{itemize}
If an overdose is suspected, call an ambulance immediately. Naloxone, a medicine that reverses opioid overdose, can be life-saving, and it is increasingly available to people at risk and their families through local programs.

\section*{Diagnosis and Investigations}
Safe prescribing involves assessing the pain, the cause, and the person's risk of harm from opioids.

\begin{itemize}
    \item \textbf{Pain and functional assessment:} the doctor asks about the cause, severity, and effect of pain on sleep, work, and activity
    \item \textbf{Medical and psychiatric history:} past dependence, mental illness, and kidney, liver, or breathing disease all affect safety
    \item \textbf{Review of all medicines:} checking for interactions, especially with alcohol, benzodiazepines, and other sedatives
    \item \textbf{Opioid risk assessment:} structured tools and questions identify people at higher risk of dependence
    \item \textbf{Drug monitoring:} in long-term or specialist care, urine tests and prescription databases (such as the local Prescription Shopping program and state real-time monitoring) are used to check adherence
    \item \textbf{Regular review:} for chronic pain, planned review visits check pain, function, side effects, and whether the dose can be reduced
\end{itemize}

\section*{Management}
Managing opioids safely means using the right dose for the right duration, planning the end of treatment from the start, and acting on problems early.

\begin{itemize}
    \item \textbf{Take exactly as prescribed:} never increase the dose, shorten the interval, or take extra tablets without advice
    \item \textbf{Short courses for acute pain:} most acute pain needs only a few days of opioids, then a step-down to simpler medicines
    \item \textbf{Plan the taper:} for longer courses, the doctor plans a gradual dose reduction so withdrawal is avoided
    \item \textbf{Prevent constipation:} drink plenty of fluids, eat fibre, stay active, and use a laxative if needed
    \item \textbf{Do not drive while sedated:} opioids can impair driving, especially in the first week, after dose increases, or with alcohol
    \item \textbf{Avoid alcohol and other sedatives:} combining them with opioids sharply increases the risk of overdose
    \item \textbf{Use the medicines management plan:} a shared written plan of doses, review dates, and when to seek help
    \item \textbf{Naloxone for at-risk people:} people on high doses or with dependence can be given naloxone and trained to use it
    \item \textbf{Opioid treatment programs:} for dependence, methadone or buprenorphine treatment provides a safe, supervised alternative
\end{itemize}

\section*{Complications}
The complications of opioids range from the unpleasant to the fatal, and most are avoidable with careful use.

\begin{itemize}
    \item \textbf{Dependence and addiction:} the brain adapts to the drug, and some people lose control of their use
    \item \textbf{Overdose and death:} respiratory depression can be fatal, especially with alcohol or other sedatives
    \item \textbf{Constipation:} persistent and sometimes severe, with a risk of bowel obstruction
    \item \textbf{Falls and fractures:} sedation and dizziness cause falls, particularly in older people
    \item \textbf{Opioid-induced hyperalgesia:} paradoxically, long-term high-dose opioids can increase sensitivity to pain
    \item \textbf{Withdrawal:} unpleasant symptoms when the medicine is reduced or stopped, managed with a supervised taper
    \item \textbf{Hormonal effects:} long-term use lowers testosterone and oestrogen, affecting energy, mood, and libido
\end{itemize}

\section*{Prognosis and Outcomes}
For acute pain, the outlook with opioids is excellent when they are used for a few days and then stopped: pain settles, the medicine is withdrawn smoothly, and the risk of harm is low. For chronic pain, the picture is more complex. Opioids provide meaningful relief for some people, but for many their benefit fades with time, and the emphasis has shifted to rehabilitation-based care that combines activity, exercise, psychological support, and, where needed, modest and stable opioid doses. This approach produces the best long-term function.

Dependence is treatable. With structured programs, medical and psychological support, and gradual withdrawal, many people regain control of their lives, and the risk of death falls markedly once opioids are reduced or replaced. The safest users are those who treat the medicine with respect: a powerful tool for severe pain, used at the right dose for the right time, kept secure, and stopped when it is no longer needed.

\section*{Prevention and Self-Care}
\begin{itemize}
    \item Ask your doctor and pharmacist about non-opioid options for pain before starting opioids
    \item Use the lowest effective dose for the shortest time, and plan from day one how the medicine will be reduced
    \item Never combine opioids with alcohol, sleeping tablets, or other sedatives
    \item Store opioids in a locked cupboard and out of reach of children and visitors
    \item Return unused tablets to a pharmacy for safe disposal; never pass them to others
    \item Keep a written list of all your medicines and show it to every doctor and pharmacist
    \item Learn the signs of overdose and how to respond, and consider naloxone training if you take high doses
\end{itemize}

\section*{Sources and Further Reading}
The following organisations provide reliable information about the safe use of opioids.

\begin{itemize}
    \item NHS. \textit{Opioids: painkillers and how to use them safely.} https://www.nhs.uk/
    \item Mayo Clinic. \textit{Opioid use and safer pain relief.} https://www.mayoclinic.org/
    \item healthdirect Australia. \textit{Opioid pain medicines: using them safely.} https://www.healthdirect.gov.au/
    \item World Health Organization (WHO). \textit{Opioid use and overdose fact sheet.} https://www.who.int/
    \item Australian Department of Health and Aged Care. \textit{Opioid medicines and safer prescribing.} https://www.health.gov.au/
    \item Australian Prescriber. \textit{Opioid analgesics: safe prescribing for pain.} https://www.nps.org.au/
\end{itemize}
\clearpage
